% sec8_7.tex — Section 8.7: Serially Correlated Observations

So far, we have assumed i.i.d.\ observations in deriving the asymptotic variance of
ML estimators. This section is concerned with ML estimation when observations
are serially correlated.

\subsection*{Two Questions}

When observations are serially correlated, there arise two distinct questions. The
first is whether the ML estimator that maximizes a likelihood function derived from
the assumption of i.i.d.\ observations is consistent and asymptotically normal when
indeed the observations are serially correlated. More precisely, suppose that the
observation vector $\mathbf{w}_t$ is ergodic stationary but not necessarily i.i.d., and consider
the ML estimator that maximizes
\begin{equation}
  Q_n(\boldsymbol{\theta})
  = \frac{1}{n}\sum_{t=1}^{n} \log f(\mathbf{w}_t;\, \boldsymbol{\theta}),
  \label{eq:8.7.1} \tag{8.7.1}
\end{equation}
where $\log f(\mathbf{w}_t;\, \boldsymbol{\theta})$ is the log likelihood for observation~$t$. This objective function
would be ($1/n$ times) the log likelihood of the sample $(\mathbf{w}_1, \ldots, \mathbf{w}_n)$ if the observations were i.i.d. Even though the log likelihood for observation~$t$,
$\log f(\mathbf{w}_t;\, \boldsymbol{\theta})$,
is correctly specified, this objective function is misspecified because it assumes,
incorrectly, that $\{\mathbf{w}_t\}$ is i.i.d. The ML estimator that maximizes this objective function is a quasi-ML estimator because the likelihood function is misspecified. Is a
quasi-ML estimator that incorrectly assumes no serial correlation consistent and
asymptotically normal?

We have actually answered this first question in the previous chapter. By
Propositions~7.5 and~7.6, the quasi-ML estimator is consistent because it is an
M-estimator, with $m(\mathbf{w}_t;\, \boldsymbol{\theta})$ given by the log likelihood for observation~$t$. For
asymptotic normality, the relevant theorem is Proposition~7.8, which spells out
conditions under which an M-estimator is asymptotically normal. Just reproducing the conclusion of Proposition~7.8, the expression for the asymptotic variance is
given by
\begin{equation}
  \avar(\hat{\boldsymbol{\theta}})
  = \bigl(\E[\mathbf{H}(\mathbf{w}_t;\, \boldsymbol{\theta}_0)]\bigr)^{-1}
    \boldsymbol{\Sigma}
    \bigl(\E[\mathbf{H}(\mathbf{w}_t;\, \boldsymbol{\theta}_0)]\bigr)^{-1},
  \label{eq:8.7.2} \tag{8.7.2}
\end{equation}
where $\boldsymbol{\theta}_0$ is the true value of $\boldsymbol{\theta}$,
$\mathbf{H}(\mathbf{w}_t;\, \boldsymbol{\theta})$ is the Hessian for observation~$t$, and
$\boldsymbol{\Sigma}$ is
the long-run variance of the score for observation~$t$ evaluated at $\boldsymbol{\theta}_0$,
$\mathbf{s}(\mathbf{w}_t;\, \boldsymbol{\theta}_0)$. This
asymptotic variance can be consistently estimated by
\begin{equation}
  \widehat{\avar}(\hat{\boldsymbol{\theta}})
  = \biggl\{\frac{1}{n}\sum_{t=1}^{n}
    \mathbf{H}(\mathbf{w}_t;\, \hat{\boldsymbol{\theta}})\biggr\}^{-1}
    \hat{\boldsymbol{\Sigma}}
    \biggl\{\frac{1}{n}\sum_{t=1}^{n}
    \mathbf{H}(\mathbf{w}_t;\, \hat{\boldsymbol{\theta}})\biggr\}^{-1},
  \label{eq:8.7.3} \tag{8.7.3}
\end{equation}
where $\hat{\boldsymbol{\Sigma}}$ is a consistent estimate of $\boldsymbol{\Sigma}$. Any of the ``nonparametric'' methods discussed in Chapter~6 (such as the VARHAC) can be used to calculate $\hat{\boldsymbol{\Sigma}}$ from the
estimated series $\{\mathbf{s}(\mathbf{w}_t;\, \hat{\boldsymbol{\theta}})\}$; there is no need to parameterize the serial correlation
in $\{\mathbf{s}(\mathbf{w}_t;\, \boldsymbol{\theta}_0)\}$.

The second question is what is the likelihood function that properly accounts
for serial correlation in observations and what are the properties of the ``genuine''
ML estimator that maximizes it. The question can be posed for conditional ML
as well as unconditional ML. In conditional ML, we divide the observation vector
$\mathbf{w}_t$ into two groups, the dependent variable $y_t$ and the regressors $\mathbf{x}_t$, and examine
the conditional distribution of $(y_1, \ldots, y_n)$ conditional on $(\mathbf{x}_1, \ldots, \mathbf{x}_n)$. Parameterizing this conditional distribution is relatively straightforward when $\{\mathbf{w}_t\}$ is i.i.d.,
because it can be written as a product over~$t$ of the conditional distribution of $y_t$
given $\mathbf{x}_t$. In fact, this i.i.d.\ case is what we have examined in this and previous
chapters. If $\{\mathbf{w}_t\}$ is serially correlated, however, the researcher does not know the
dynamic interaction between $y_t$ and $\mathbf{x}_t$ well enough to be able to write down the
conditional distribution.\footnote{For a discussion of conditional ML when the complete specification of the dynamic interaction is not possible or undesirable, see Wooldridge (1994, Section~5). For a discussion of possible restrictions on the dynamic
interaction (such as ``Granger causality'' and ``weak exogeneity''), see, e.g., Davidson and MacKinnon (1993,
Section~18.2).}
In the rest of this section, we deal only with examples
of unconditional ML with serially correlated observations.

\subsection*{Unconditional ML for Dependent Observations}

We start out with the simplest case of a univariate series. Let $(y_0, y_1, \ldots, y_n)$ be the
sample. (For notational convenience, we assume that the sample includes observation for $t = 0$.) When the observations are serially correlated, the density function
of the sample can no longer be expressed as a product of individual densities. However, it can be written down in terms of a series of conditional probability density
functions. The joint density function of $(y_0, y_1)$ can be written as
\begin{equation}
  f(y_0, y_1) = f(y_1 \mid y_0)\,f(y_0).
  \label{eq:8.7.4} \tag{8.7.4}
\end{equation}
Similarly, for $(y_0, y_1, y_2)$,
\begin{equation}
  f(y_0, y_1, y_2) = f(y_2 \mid y_1, y_0)\,f(y_0, y_1).
  \label{eq:8.7.5} \tag{8.7.5}
\end{equation}
Substituting \eqref{eq:8.7.4} into this, we obtain
\begin{equation}
  f(y_0, y_1, y_2) = f(y_2 \mid y_1, y_0)\,f(y_1 \mid y_0)\,f(y_0).
  \label{eq:8.7.6} \tag{8.7.6}
\end{equation}
A repeated use of this sort of sequential substitution produces
\begin{equation}
  f(y_0, y_1, \ldots, y_n)
  = \prod_{t=1}^{n} f(y_t \mid y_{t-1}, \ldots, y_1, y_0)\,f(y_0).
  \label{eq:8.7.7} \tag{8.7.7}
\end{equation}
If the conditional density $f(y_t \mid y_{t-1}, \ldots, y_1, y_0)$ and the unconditional density
$f(y_0)$ are parameterized by a finite-dimensional parameter vector $\boldsymbol{\theta}$, then ($1/n$
times) the log likelihood of the sample $(y_0, y_1, \ldots, y_n)$ can be written as
\begin{equation}
  \frac{1}{n}\sum_{t=1}^{n}
    \log f(y_t \mid y_{t-1}, \ldots, y_1, y_0;\, \boldsymbol{\theta})
  + \frac{1}{n}\log f(y_0;\, \boldsymbol{\theta}).
  \label{eq:8.7.8} \tag{8.7.8}
\end{equation}
The ML estimator $\tilde{\boldsymbol{\theta}}$ of the true parameter value $\boldsymbol{\theta}_0$ is the $\boldsymbol{\theta}$ that maximizes this log
likelihood function. We will call this ML estimator the \textbf{exact ML estimator}.

\subsection*{ML Estimation of AR(1) Processes}

In Section~6.2, we considered in some detail first-order autoregressive (AR(1))
processes. If we assume for an AR(1) process that the error term is normally
distributed, we obtain a \textbf{Gaussian AR(1) process}:
\begin{equation}
  y_t = c_0 + \phi_0 y_{t-1} + \varepsilon_t,
  \label{eq:8.7.9} \tag{8.7.9}
\end{equation}
with $\varepsilon_t \sim \text{i.i.d.\ } N(0, \sigma_0^2)$. We assume the stationarity condition that $|\phi_0| < 1$.
So $\{y_t\}$ is ergodic stationary by Proposition~6.1(d). From Section~6.2, we know
that the unconditional mean of $y_t$ is $c_0/(1 - \phi_0)$ and the unconditional variance
is $\sigma_0^2/(1 - \phi_0^2)$. Furthermore, since $y_t$ can be written as a weighted average of
$(\varepsilon_t, \varepsilon_{t-1}, \ldots)$, its distribution is normal when $\varepsilon_t$ is normal. Therefore,
\begin{equation}
  y_0 \sim N\!\biggl(\frac{c_0}{1 - \phi_0},\;
    \frac{\sigma_0^2}{1 - \phi_0^2}\biggr).
  \label{eq:8.7.10} \tag{8.7.10}
\end{equation}
Also, it follows directly from \eqref{eq:8.7.9} and the normality assumption that
\begin{equation}
  y_t \mid (y_{t-1}, \ldots, y_1, y_0)
  \sim N(c_0 + \phi_0 y_{t-1},\, \sigma_0^2).
  \label{eq:8.7.11} \tag{8.7.11}
\end{equation}
Therefore, using the formula \eqref{eq:8.7.7}, the joint density of $(y_0, y_1, \ldots, y_n)$ for a
hypothetical parameter value $\boldsymbol{\theta} \equiv (c, \phi, \sigma^2)'$ can be written as
\begin{multline}
  f(y_0, y_1, \ldots, y_n;\, \boldsymbol{\theta})
  = \prod_{t=1}^{n}\biggl\{
    \frac{1}{\sqrt{2\pi\sigma^2}}
    \exp\biggl[-\frac{1}{2\sigma^2}(y_t - c - \phi y_{t-1})^2\biggr]\biggr\} \\
  \times \frac{1}{\sqrt{2\pi\frac{\sigma^2}{1-\phi^2}}}
    \exp\Biggl[-\frac{\Bigl(y_0 - \frac{c}{1-\phi}\Bigr)^2}
    {2\frac{\sigma^2}{1-\phi^2}}\Biggr].
  \label{eq:8.7.12} \tag{8.7.12}
\end{multline}
Taking logs and dividing both sides by~$n$, we obtain the objective function for exact
ML:
\begin{multline}
  \tilde{Q}_n(\boldsymbol{\theta})
  = \frac{1}{n}\sum_{t=1}^{n}
    \biggl\{-\frac{1}{2}\log(2\pi) - \frac{1}{2}\log(\sigma^2)
    - \frac{1}{2\sigma^2}(y_t - c - \phi y_{t-1})^2\biggr\} \\
  + \frac{1}{n}\biggl\{-\frac{1}{2}\log(2\pi)
    - \frac{1}{2}\log\!\Bigl(\frac{\sigma^2}{1-\phi^2}\Bigr)
    - \frac{\Bigl(y_0 - \frac{c}{1-\phi}\Bigr)^2}
      {2\frac{\sigma^2}{1-\phi^2}}\biggr\}.
  \label{eq:8.7.13} \tag{8.7.13}
\end{multline}
The exact ML estimator $\tilde{\boldsymbol{\theta}}$ of the AR(1) process is the $\boldsymbol{\theta}$ that maximizes this objective function. Assuming an interior solution, the estimator solves the first-order
conditions (the likelihood equations) obtained by setting $\frac{\partial \tilde{Q}_n}{\partial \boldsymbol{\theta}}$
to zero. They are a
system of equations that are nonlinear in $\boldsymbol{\theta}$, so finding $\tilde{\boldsymbol{\theta}}$ requires iterative algorithms such as those described in Section~7.5.

\subsection*{Conditional ML Estimation of AR(1) Processes}

The source of the nonlinearity in the likelihood equations is the log likelihood of
$y_0$. As an alternative, consider dropping this term and maximizing
\begin{equation}
  Q_n(\boldsymbol{\theta})
  = \frac{1}{n}\sum_{t=1}^{n}
    \biggl\{-\frac{1}{2}\log(2\pi) - \frac{1}{2}\log(\sigma^2)
    - \frac{1}{2\sigma^2}(y_t - c - \phi y_{t-1})^2\biggr\}.
  \label{eq:8.7.14} \tag{8.7.14}
\end{equation}
Using the relation $f(y_0, y_1, \ldots, y_n) = f(y_1, \ldots, y_n \mid y_0)\,f(y_0)$ and \eqref{eq:8.7.12}, we
see that this objective function is ($1/n$ times) the log of the likelihood conditional
on the first observation $y_0$,
\begin{equation}
  f(y_1, y_2, \ldots, y_n \mid y_0;\, \boldsymbol{\theta})
  = \prod_{t=1}^{n}\biggl\{
    \frac{1}{\sqrt{2\pi\sigma^2}}
    \exp\biggl[-\frac{1}{2\sigma^2}(y_t - c - \phi y_{t-1})^2\biggr]\biggr\}.
  \label{eq:8.7.15} \tag{8.7.15}
\end{equation}

Let $\hat{\boldsymbol{\theta}}$ be the $\boldsymbol{\theta}$ that maximizes the log conditional likelihood \eqref{eq:8.7.14} (conditional
on~$y_0$). This ML estimator will henceforth be called the \textbf{$y_0$-conditional ML estimator}. Here, conditioning is on $y_0$, which makes the estimator here conceptually
different from the conditional ML estimators of the previous sections where conditioning is on $\mathbf{x}_t$.

Before examining the relationship between the two ML estimators, the exact
ML estimator $\tilde{\boldsymbol{\theta}}$ and the $y_0$-conditional ML estimator $\hat{\boldsymbol{\theta}}$, we give two derivations of
the asymptotic properties of the latter. The first derivation is based on the explicit
solution for the estimator. Recall from Section~1.5 and Example~7.2 that the objective function in the ML estimation of the linear regression model $y_t = \mathbf{x}_t'\boldsymbol{\beta}_0 + \varepsilon_t$
for i.i.d.\ observations is
\begin{equation}
  \frac{1}{n}\sum_{t=1}^{n}
    \biggl\{-\frac{1}{2}\log(2\pi) - \frac{1}{2}\log(\sigma^2)
    - \frac{1}{2\sigma^2}(y_t - \mathbf{x}_t'\boldsymbol{\beta})^2\biggr\}.
  \label{eq:8.7.16} \tag{8.7.16}
\end{equation}
It was shown in Section~1.5, and is very easy to show, that the ML estimator of
$\boldsymbol{\beta}_0$ is the OLS coefficient estimator and the ML estimator of $\sigma_0^2$ equals the sum
of squared residuals divided by the sample size. Now, with $\mathbf{x}_t = (1, y_{t-1})'$ and
$\boldsymbol{\beta}_0 = (c_0, \phi_0)'$, the objective function for the linear regression model reduces to
$Q_n(\boldsymbol{\theta})$ in \eqref{eq:8.7.14}, the conditional log likelihood for the AR(1) process. Therefore,
algebraically, the conditional ML estimator $(\hat{c}, \hat{\phi})$ of $(c_0, \phi_0)$ is numerically the
same as the OLS coefficient estimator in the regression of $y_t$ on a constant and $y_{t-1}$,
and the conditional ML estimator $\hat{\sigma}^2$ of $\sigma_0^2$ is the sum of squared residuals divided
by~$n$. Regarding the asymptotic properties of the estimator, we have shown in
Section~6.4 that $(\hat{c}, \hat{\phi})$ is consistent and asymptotically normal with the asymptotic
variance indicated in Proposition~6.7 for $p = 1$ and that $\hat{\sigma}^2$ is consistent for $\sigma_0^2$.
Therefore, for the asymptotic normality of the conditional ML estimator $(\hat{c}, \hat{\phi})$,
which maximizes the $y_0$-conditional likelihood for normal errors, the normality
assumption is not needed.

The second derivation is more indirect but nevertheless more instructive. As
just noted, the objective function $Q_n(\boldsymbol{\theta})$ coincides with \eqref{eq:8.7.16} with
$\mathbf{x}_t = (1, y_{t-1})'$
and $\boldsymbol{\beta}_0 = (c_0, \phi_0)'$. Since $\{y_t, \mathbf{x}_t\}$ is merely ergodic stationary and not necessarily
i.i.d., the conditional ML estimator $(\hat{c}, \hat{\phi}, \hat{\sigma}^2)$ can be viewed as the quasi-ML estimator considered in Proposition~7.6. It has been verified in Example~7.8 that the
objective function satisfies all the conditions of Proposition~7.6 if $\E(\mathbf{x}_t\mathbf{x}_t')$ where
$\mathbf{x}_t = (1, y_{t-1})'$ is nonsingular. We have shown in Section~6.4 that this nonsingularity condition is satisfied for AR(1) processes if $\sigma_0^2 > 0$. Thus, the conditional ML
is consistent. To prove asymptotic normality, view the conditional ML estimator
as an M-estimator with $\mathbf{w}_t = (y_t, y_{t-1})'$ and
\begin{equation}
  m(\mathbf{w}_t;\, \boldsymbol{\theta})
  = -\frac{1}{2}\log(2\pi) - \frac{1}{2}\log(\sigma^2)
    - \frac{1}{2\sigma^2}(y_t - c - \phi y_{t-1})^2.
  \label{eq:8.7.17} \tag{8.7.17}
\end{equation}
The relevant theorem, therefore, is Proposition~7.8. Let $\mathbf{s}(\mathbf{w}_t;\, \boldsymbol{\theta})$ be the score for
observation~$t$ and let $\mathbf{H}(\mathbf{w}_t;\, \boldsymbol{\theta})$ be the Hessian for observation~$t$ associated with this
$m$ function, and consider the conditions (1)--(5) of Proposition~7.8. The discussion
in Example~7.10 implies that all the conditions except for~(3) hold if $\E(\mathbf{x}_t\mathbf{x}_t')$ where
$\mathbf{x}_t = (1, y_{t-1})'$ is nonsingular, which as just noted is satisfied if $\sigma_0^2 > 0$. This leaves
condition~(3) that
$\frac{1}{\sqrt{n}}\sum_{t=1}^{n} \mathbf{s}(\mathbf{w}_t;\, \boldsymbol{\theta}_0) \dto N(\mathbf{0}, \boldsymbol{\Sigma})$.
In the AR(1) model, $f(y_t \mid y_{t-1}, \ldots, y_1, y_0) = f(y_t \mid y_{t-1})$. As you will show in Review Question~1, this
special feature of the joint density function implies that the sequence
$\{\mathbf{s}(\mathbf{w}_t;\, \boldsymbol{\theta}_0)\}$ is
a martingale difference sequence. Since $\mathbf{s}(\mathbf{w}_t;\, \boldsymbol{\theta}_0)$ is a function of
$\mathbf{w}_t = (y_t, y_{t-1})'$
and $\{y_t\}$ is ergodic stationary, the sequence $\{\mathbf{s}(\mathbf{w}_t;\, \boldsymbol{\theta}_0)\}$ also is ergodic stationary.
So by the ergodic stationary martingale difference CLT of Chapter~2, we can claim
that condition~(3) is satisfied with
\begin{equation}
  \boldsymbol{\Sigma}
  = \E\bigl[\mathbf{s}(\mathbf{w}_t;\, \boldsymbol{\theta}_0)\,
    \mathbf{s}(\mathbf{w}_t;\, \boldsymbol{\theta}_0)'\bigr].
  \label{eq:8.7.18} \tag{8.7.18}
\end{equation}

Finally, it is easy to show the information matrix equality for observation~$t$ that
$-\E[\mathbf{H}(\mathbf{w}_t;\, \boldsymbol{\theta}_0)]
= \E[\mathbf{s}(\mathbf{w}_t;\, \boldsymbol{\theta}_0)\,\mathbf{s}(\mathbf{w}_t;\, \boldsymbol{\theta}_0)']$
(a proof is in Example~7.10). Substituting this and \eqref{eq:8.7.18} into \eqref{eq:8.7.2}, we conclude that the asymptotic variance of the
$y_0$-conditional ML estimator $\hat{\boldsymbol{\theta}}$ is given by
\begin{equation}
  \avar(\hat{\boldsymbol{\theta}})
  = -\bigl(\E[\mathbf{H}(\mathbf{w}_t;\, \boldsymbol{\theta}_0)]\bigr)^{-1}
  = \bigl(\E[\mathbf{s}(\mathbf{w}_t;\, \boldsymbol{\theta}_0)\,
    \mathbf{s}(\mathbf{w}_t;\, \boldsymbol{\theta}_0)']\bigr)^{-1}.
  \label{eq:8.7.19} \tag{8.7.19}
\end{equation}
That is, the usual conclusion for a correctly specified ML estimator for i.i.d.\ observations is applicable to the $y_0$-conditional ML estimator that maximizes the correctly specified conditional likelihood \eqref{eq:8.7.14} for dependent observations.\footnote{This argument can be applied to models more general than autoregressive processes. See Lemma~5.2 of
Wooldridge (1994).}

The difference between \eqref{eq:8.7.13} and \eqref{eq:8.7.14} is the log likelihood of $y_0$ divided
by~$n$. If the sample size $n$ is sufficiently large, then the first observation $y_0$ makes a
negligible contribution to the likelihood of the sample. The exact ML estimator and
the $y_0$-conditional ML estimator turn out to have the same asymptotic distribution,
provided that $|\phi| < 1$ (see, e.g., Fuller, 1996, Sections~8.1 and~8.4). For this
reason, in most applications the parameters of an autoregression are estimated by
OLS ($y_0$-conditional ML) rather than exact ML.

\subsection*{Conditional ML Estimation of AR(\texorpdfstring{$p$}{p}) and VAR(\texorpdfstring{$p$}{p}) Processes}

A \textbf{Gaussian AR($p$)} process can be written as
\begin{equation}
  y_t = c_0 + \phi_{01} y_{t-1} + \phi_{02} y_{t-2} + \cdots + \phi_{0p} y_{t-p}
  + \varepsilon_t,
  \label{eq:8.7.20} \tag{8.7.20}
\end{equation}
with $\varepsilon_t \sim \text{i.i.d.\ } N(0, \sigma_0^2)$. The conditional ML estimation of the AR($p$) process is
completely analogous to the AR(1) case. If the sample is $(y_{-p+1}, y_{-p+2}, \ldots, y_0,
y_1, \ldots, y_n)$, then ($1/n$ times) the log likelihood of the sample conditional on
$(y_{-p+1}, \ldots, y_0)$ is \eqref{eq:8.7.16} with
\[
  \mathbf{x}_t = (1, y_{t-1}, \ldots, y_{t-p})' \quad\text{and}\quad
  \boldsymbol{\beta}_0 = (c_0, \phi_{01}, \phi_{02}, \ldots, \phi_{0p})'.
\]
Therefore, the conditional ML (conditional on $(y_{-p+1}, \ldots, y_0)$) estimator
of $(c_0, \phi_{01}, \ldots, \phi_{0p})$ is numerically the same as the OLS coefficient estimator in
the regression of $y_t$ on a constant and $p$ lagged values of $y_t$. The conditional ML
estimate of $\sigma_0^2$ is the sum of squared residuals divided by the sample size. By
Proposition~6.7, the conditional ML estimator of the coefficients is consistent and
asymptotically normal if the stationarity condition (that the roots of
$1 - \phi_{01}z - \cdots - \phi_{0p}z^p = 0$ be greater than~1 in absolute value) is satisfied. Again as in the AR(1)
case, the error term does not need to be normal for consistency and asymptotic
normality. The exact ML estimates and the conditional ML estimates have the
same asymptotic distribution if the stationarity condition is satisfied.

These results generalize readily to vector processes. A \textbf{Gaussian VAR($p$)}
($p$-th order vector autoregression) can be written as
\begin{equation}
  \underset{(M \times 1)}{\mathbf{y}_t}
  = \mathbf{c}_0 + \boldsymbol{\Phi}_{01}\mathbf{y}_{t-1}
  + \cdots + \boldsymbol{\Phi}_{0p}\mathbf{y}_{t-p}
  + \boldsymbol{\varepsilon}_t,
  \label{eq:8.7.21} \tag{8.7.21}
\end{equation}
with $\boldsymbol{\varepsilon}_t \sim \text{i.i.d.\ } N(\mathbf{0}, \boldsymbol{\Omega}_0)$.
The objective function in the conditional ML estimation
is ($1/n$ times) the log conditional likelihood:
\begin{equation}
  Q_n(\boldsymbol{\theta})
  = -\frac{M}{2}\log(2\pi)
    + \frac{1}{2}\log(|\boldsymbol{\Omega}^{-1}|)
    - \frac{1}{2n}\sum_{t=1}^{n}
      \{(\mathbf{y}_t - \boldsymbol{\Pi}'\mathbf{x}_t)'
      \boldsymbol{\Omega}^{-1}
      (\mathbf{y}_t - \boldsymbol{\Pi}'\mathbf{x}_t)\},
  \label{eq:8.7.22} \tag{8.7.22}
\end{equation}
where
\begin{equation}
  \boldsymbol{\Pi}' \equiv
  \begin{bmatrix}
    \mathbf{c} & \boldsymbol{\Phi}_1 & \boldsymbol{\Phi}_2 & \cdots & \boldsymbol{\Phi}_p
  \end{bmatrix}, \quad
  \mathbf{x}_t \equiv
  \begin{bmatrix} 1 \\ \mathbf{y}_{t-1} \\ \mathbf{y}_{t-2} \\ \vdots \\ \mathbf{y}_{t-p} \end{bmatrix}.
  \label{eq:8.7.23} \tag{8.7.23}
\end{equation}
This coincides with the objective function in the ML estimation of the multivariate regression model considered in Section~8.4. Therefore, the conditional ML
estimate of $(\mathbf{c}_0, \boldsymbol{\Pi}_0)$ is numerically the same as the equation-by-equation OLS estimator. It is consistent and asymptotically normal if the coefficient matrices satisfy
the stationarity condition (that all the roots of
$|\mathbf{I}_M - \boldsymbol{\Phi}_{01}z - \cdots - \boldsymbol{\Phi}_{0p}z^p| = 0$ be
greater than~1 in absolute value). This result does not require the error vector $\boldsymbol{\varepsilon}_t$ to
be normally distributed.

\bigskip
\noindent\rule{\textwidth}{0.4pt}
\textbf{QUESTIONS FOR REVIEW}
\medskip

\begin{enumerate}
\item (Score is m.d.s.)\quad
  Consider the Gaussian AR(1) process \eqref{eq:8.7.9} and let
  $\mathbf{s}(\mathbf{w}_t;\, \boldsymbol{\theta})$
  be the score for observation~$t$ (the gradient of \eqref{eq:8.7.17}) with
  $\mathbf{w}_t = (y_t, y_{t-1})'$
  and $\boldsymbol{\theta} = (c, \phi, \sigma^2)'$.

  \begin{enumerate}[label=(\alph*)]
  \item Show that
    \[
      \mathbf{s}(\mathbf{w}_t;\, \boldsymbol{\theta})
      = \begin{bmatrix}
          \frac{1}{\sigma^2}\mathbf{x}_t \cdot (y_t - \mathbf{x}_t'\boldsymbol{\beta}) \\[6pt]
          -\frac{1}{2\sigma^2} + \frac{1}{2\sigma^4}(y_t - \mathbf{x}_t'\boldsymbol{\beta})^2
        \end{bmatrix},
    \]
    where $\boldsymbol{\beta} = (c, \phi)'$ and $\mathbf{x}_t = (1, y_{t-1})'$.
  \end{enumerate}
\end{enumerate}
